\documentclass[10pt,a4paper]{amsart}
\usepackage[margin=19mm]{geometry}
\usepackage{amsmath,amssymb,mathtools}
\usepackage[scr=rsfs,cal=euler]{mathalfa}
\usepackage{xfrac}
\usepackage[unicode,hidelinks,bookmarks=false]{hyperref}
\newcommand{\ie}{i.e.\ }
\newcommand{\cf}{cf.\ }
\newcommand{\resp}{respectively\ }
\newcommand{\Subsection}{\subsection}
\providecommand{\nobreakdash}{\nobreak\hbox{-}}
\setlength{\emergencystretch}{2em}
\multlinegap0pt
\usepackage{amssymb}
\usepackage{bm}
\newtheorem{lem}{Lemma}
\def\Z{{\mathbin{\mathbb Z}}}
\def\al{\alpha}
\def\ep{\epsilon}
\def\ge{\geqslant}
\def\iy{\infty}
\def\l{\label}
\def\om{\omega}
\def\sb{\subseteq}
\def\ta{\tau}
\def\th{\theta}
\title{Generalized Thomas--Yau Uniqueness Theorems}
\author{Yohsuke Imagi}
\date{Source: arXiv:2012.14647v2 (CC BY 4.0)}
\begin{document}
\maketitle

\noindent\emph{Source and licence.} Generalized Thomas--Yau Uniqueness Theorems. Version arXiv:2012.14647v2 (CC BY 4.0), distributed by arXiv under the Creative Commons Attribution 4.0 International licence, http://creativecommons.org/licenses/by/4.0/, as recorded in the arXiv licence metadata for this exact version and shown on the abstract page https://arxiv.org/abs/2012.14647v2. Full licence notice: \texttt{licenses/arxiv-2012.14647-CC-BY-4.0.txt}.\par\smallskip

\noindent\textit{Editorial scope.} Counted unit: the article's lem lem (source0.tex lines 555--563). The article's complete original proof of it is delivered with it (source0.tex lines 564--581, one \texttt{proof} environment of 18 lines). No mathematics is written, repaired or re-derived: the statement and the proof are the article's own lines, in the article's own notation, and no retained line is substituted or re-flowed.  No further source-local context is needed: every cross-reference of every retained line resolves inside the two delivered spans.  Nothing was omitted from any retained span, so the package carries no omission record.  No retained line contains a citation, so the wrapper carries no bibliography and no bibliography entry.  No retained line contains a figure environment, an image-inclusion command or drawing code, so no image is delivered and the package declares no asset dependency.  Editorial formatting only, in addition: an AMS \texttt{amsart} preamble that restates the article's own declarations and macros used by the retained lines, namely the environments \texttt{lem} on separate counters, together with the standard packages those lines need.  The retained environments print in this document as Lemma 1 in order of appearance, whereas the article prints the same items with the numbering of its own full document.  The complete native source of the article as distributed in the arXiv e-print for this exact version is bundled unchanged as \texttt{sources/106377/source0.tex}.
\begin{lem}\l{P0}
Let $(X,\om)$ be a symplectic manifold of real dimension $2n,$ $\Z_k$-graded with $k\ge n,$
equipped with a compatible almost complex-structure $J$ and equipped with a Hermitian metric $g.$
Let $L_1,L_2\sb X$ be two closed immersed $\Z_k$-graded Lagrangians with the same grading on $L_1\cap L_2.$
Denote by $S\sb L_1\cap L_2$ the set of points at which $L_1,L_2$ have at least one common tangent space.
Then for every open neighbourhood $U\sb X$ of $S$
there exists $\ep>0$ such that every Hamiltonian $\ep$-perturbation of $L_2$ that intersects $L_1$ generically $($that is, only at transverse double points$)$
has no intersection point with $L_1\cap U$ of index $0$ or $n$ modulo $k\Z.$
\end{lem}

\begin{proof}
If this fails there are an open neighbourhood $U\sb X$ of $S,$
a sequence of Hamiltonian perturbations $L_2^i$ of $L_2$ all transverse to $L_1$ and tending smoothly to $L_2,$
and a sequence of intersection points $x^i\in L_1\cap L_2^i\-U$ of index $0$ or $n.$
By hypothesis $L_1\cap L_2^i\-U$ is compact, so there is a subsequence of $x^i$ tending to some point $x\in L_1\cap L_2\-U.$
Suppose now that $k<\iy;$
the other case $k=\iy$ may be treated in the same manner.
Give $L_2^i$ a grading $\al^i$ by the obvious homotopy,
and write $\al^i=\exp(\sqrt{-1}\ta^i)\al$ over $x^i.$
By hypothesis $L_1,L_2$ have the same grading at $x^i$ so we may suppose that for each $i$ we have $|\ta^i|<2\pi/k$ and the sequence $\ta^i$ tends to $0$ as $i\to\iy.$
Denote by $\th_1^i,\cdots,\th_n^i\in(0,\pi)$ the $n$ angles between the two tangent spaces at $x^i.$
Then
$2\th_1^i+\cdots+2\th_n^i+k\ta^i=0\text{ or }n\pi\text{ modulo }2k\pi\Z.$
In fact, since $k\ge n$ it follows that this holds without modulo $2k\pi\Z.$
So the limits of the $n$ angles sum up to $0$ or $n\pi;$
that is, at the limit point $x$ there is a common tangent space to $L_1,L_2.$
But this contradicts the definition of $S.$
\end{proof}

\end{document}
